\chapter*{Preface}
\addcontentsline{toc}{chapter}{Preface}

Spectroscopy turns small departures from a smooth continuum into measurements
of temperature, density, composition, radiation field, motion, and geometry.
That power comes with a practical difficulty: the meaning of a line is never
contained in its wavelength alone. A transition becomes useful only when its
atomic or molecular physics, excitation mechanism, optical depth, environment,
instrumental window, and degeneracies are considered together.

\emph{Line Atlas} is designed as a bridge between a line list and a textbook.
Each chapter develops the reasoning behind common measurements, while the
reference atlas provides a compact route back to individual transitions. The
companion website uses the same catalogue so that wavelength and redshift
questions can be explored interactively. Values in the atlas are starting
points, not substitutes for the critically evaluated source data and modelling
appropriate to a particular analysis.

This first edition emphasizes extragalactic work and the baryon cycle: H~II
regions, the neutral and molecular interstellar medium, stellar feedback, the
circumgalactic and intergalactic media, and active galactic nuclei. Radio,
infrared, ultraviolet, optical, and X-ray diagnostics are treated as connected
views of the same evolving gas.
